\paragraph{C.3m (Refuted, with an explicit counterexample): ordinary $J^1$ surjectivity fails in every dimension at least four.}
For every $N\ge4$ we construct an excellent normal
strictly henselian local complex ring $R$ of dimension
$N$, with regular puncture $X$ and a regular punctured
principal divisor $Y$, such that
\begin{equation}\label{eq:C3-ordinary-nonsurjection}
 \operatorname{Pic}(X)_{\mathrm{alg}}=0,
 \qquad
 \operatorname{Pic}^0(C)(\mathbb C)
       \lhook\joinrel\longrightarrow
       \operatorname{Pic}(Y)_{\mathrm{alg}},
\end{equation}
where $C$ is the explicitly specified Fermat plane
quartic. The right-hand injection has nontorsion
image. Thus ordinary $J^1(X)\to J^1(Y)$ is not
surjective, even with rational coefficients.
All algebraic equivalences here use families on the
fixed scheme times a smooth connected parameter
curve. The argument below works with actual line
bundles throughout.

\emph{1. The algebraic contraction and its divisor.}
Put $m=N-2\ge2$, and define
\[
\begin{aligned}
 F(z)&=z_0^4+z_1^4+z_2^4, & C&=V(F)\subset\mathbb P^2_{\mathbb C},\\
 H(z)&=z_0^2z_1^2+z_0^2z_2^2+z_1^2z_2^2,&
 \ell(x,z)&=x_0z_0+x_1z_1+x_2z_2.
\end{aligned}
\]
Here $x_0,\ldots,x_m$ are homogeneous coordinates on
$\mathbb P^m$. Set
\[
 B=\mathbb P^m_x\times\mathbb P^2_z,
 \qquad \mathcal N=\mathcal O_B(-1,-1),
 \qquad T=\operatorname{Tot}(\mathcal N).
\]
If $t$ denotes the tautological fibre vector, the
expression
\[
                         F+t\ell H
\]
is a well-defined section of the pullback of
$\mathcal O_B(0,4)$ to $T$. Define $Z$ as its zero
scheme. Its zero-section divisor is
\[
 V=\mathbb P^m\times C,
 \qquad
 \mathcal N|_V=L^{-1},
 \qquad L=\mathcal O_{\mathbb P^m}(1)\boxtimes K_C.
\]
The curve $C$ is smooth, and $Z$ is smooth in a
neighbourhood of $V$, since its differential there
has the nonzero $dF$ component. The projection of
$Z$ onto $B$ is the graph
$t=-F/(\ell H)$ on the dense open $\ell H\ne0$.
The locus $F=\ell H=0$ has codimension two in $B$,
so its inverse image cannot be another component
of $Z$. Thus $Z$ is integral.

The standard Segre cone contraction
\[
      T\longrightarrow\mathbb A^{3(m+1)},
           \qquad w_{ij}=t x_i z_j
\]
is proper, contracts the zero section, and is an
isomorphism outside it. Let $W$ be the reduced
closed image of $Z$, and write $v$ for the vertex.
This gives a completely algebraic finite-type
construction. Explicit equations for $W$ are the
radical of the following ideal: all $2$ by $2$
minors of $(w_{ij})$, together with, for every
ordered quadruple $I=(i_1,i_2,i_3,i_4)$,
\begin{equation}\label{eq:C3-deformed-cone-equations}
 \sum_{j=0}^{2}\prod_{r=1}^{4}w_{i_rj}
 +f\sum_{0\le a<b\le2}
       w_{i_1a}w_{i_2a}w_{i_3b}w_{i_4b}=0,
 \qquad f=w_{00}+w_{11}+w_{22}.
\end{equation}
Indeed on the rank-one locus the left side is
$t^4x_{i_1}\cdots x_{i_4}(F+t\ell H)$;
choosing a nonzero row proves the converse away
from $v$. The vertex lies in the closure of this
integral open, so the reduced image is integral.

Let $W^{\mathrm n}$ be the normalization of the
original reduced affine scheme $W$; it is finite.
For the proof of its local properties choose an
open neighbourhood $W^\circ$ of $v$ such that
$Z^\circ=Z\times_W W^\circ$ is smooth.
Such a neighbourhood exists: the image of the
complement of the smooth locus is closed by
properness and does not contain $v$.
The normal smooth scheme $Z^\circ$ factors
through $(W^{\mathrm n})^\circ$, the inverse image
of $W^\circ$; the factor is proper, birational
and an isomorphism away from the vertex.
There is only one point $v^{\mathrm n}$ above $v$.
In fact the factorization is surjective, and the
connected fibre $V$ cannot map onto more than one
point in the finite fibre of the normalization.
Its residue field is $\mathbb C$. Define
\[
 A=\mathcal O_{W^{\mathrm n},v^{\mathrm n}},
 \qquad R=A^{\mathrm{sh}},
 \qquad
 \mathcal Z=Z^\circ\times_{(W^{\mathrm n})^\circ}
                         \operatorname{Spec}R.
\]
In particular $R$ is specified by the original
polynomial equations and normalization, followed
by localization at its unique vertex and strict
henselization. The auxiliary shrinking only
proves regularity and introduces no choice in $R$.
Then $R$ has all the asserted excellence, normality
and dimension properties. The regular scheme
$\mathcal Z$ is proper over $R$, its exceptional
fibre is the irreducible smooth divisor $V$, and
its complement is the regular puncture $X$.
These statements are preserved under the pointed
etale limit defining the strict henselization.

The displayed $f$ is a nonzero nonunit in $R$, and
on $Z$ it equals $t\ell$. The global equation of
the strict transform of its divisor is
\[
 D=V(\ell,F)\subset T,
\]
the total space of $L_H^{-1}$ over
\[
 H_0=V(\ell)\subset V,
 \qquad L_H=L|_{H_0}.
\]
Indeed setting $\ell=0$ in the equation of $Z$
leaves $F=0$. The surjection
\[
 \mathcal O_C^{m+1}\longrightarrow K_C,
 \qquad (a_i)\longmapsto a_0z_0+a_1z_1+a_2z_2
\]
identifies $H_0$ with the projective bundle of
lines in its rank-$m$ kernel. In particular
$H_0$ and $D$ are smooth. Set $D^\circ=D\cap Z^\circ$;
its morphism to $(W^{\mathrm n})^\circ$ is the
restriction of the factorization already constructed.
The base change
\[
 \mathcal D=D^\circ\times_{(W^{\mathrm n})^\circ}
                              \operatorname{Spec}R
\]
is regular and proper over $R$. Its exceptional
divisor is $H_0$, with normal bundle $L_H^{-1}$,
and
\[
            \mathcal D\setminus H_0=Y=X\cap V(f).
\]
This proves regularity of the required punctured
divisor, without asserting any projective-normality
property of its section ring.

\emph{2. A full-rank first-neighbourhood deformation.}
Let $V_2\subset\mathcal Z$ be defined by the square
of the ideal of $V$. It is a square-zero thickening
of $V$ by $L$. Its class
\[
                   \kappa\in H^1(V,T_V\otimes L)
\]
is obtained by taking the boundary of $\ell H$
in the normal sequence of $V\subset B$, twisted
by $L$. Its curve component is
\begin{equation}\label{eq:C3-KS-components}
       \kappa=\sum_{i=0}^{2}x_i\otimes\delta(z_iH),
 \qquad
 \delta:H^0(C,\mathcal O_C(5))\longrightarrow H^1(C,\mathcal O_C),
\end{equation}
where $\delta$ is the boundary of
\[
 0\longrightarrow T_C\otimes K_C=\mathcal O_C
 \longrightarrow T_{\mathbb P^2}|_C\otimes K_C
 \longrightarrow\mathcal O_C(5)\longrightarrow0.
\]
To verify the class directly, choose local vector
fields lifting the normal displacement $\ell H$.
Their differences on overlaps are the derivations
which identify the local retractions of $V_2$
onto $V$; this is precisely the boundary cocycle.
The sign depends on the choice of retraction and
does not affect the injectivity below.

The Euler sequence twisted by $K_C=\mathcal O_C(1)$
and Serre duality give
$H^1(C,T_{\mathbb P^2}|_C\otimes K_C)=0$.
Thus $\delta$ is surjective. More explicitly its
cokernel description is the degree-five part of
the Fermat Jacobian ring:
\[
 H^1(C,\mathcal O_C)
   \simeq
 \left(\frac{\mathbb C[z_0,z_1,z_2]}
                  {(z_0^3,z_1^3,z_2^3)}\right)_5.
\]
Here the image of quadratic vector fields consists
of the degree-five Jacobian ideal, modulo the
three multiples $z_iF$ already zero on $C$.
It has dimension $15$, while
$h^0(\mathcal O_C(5))=18$; the remaining global
tangent section is the kernel
$H^0(T_C\otimes K_C)$, so introduces no additional
relation. The three images of $z_iH$ are
\[
             z_0z_1^2z_2^2,\qquad
             z_0^2z_1z_2^2,\qquad
             z_0^2z_1^2z_2.
\]
They form a basis of that degree-five part.

Contraction with $\kappa$ consequently gives an
injective linear map
\begin{equation}\label{eq:C3-KS-full-rank}
 c_\kappa:H^0(C,K_C)
       =H^0(V,\Omega_V^1)
       \longrightarrow H^1(V,L).
\end{equation}
Indeed $H^1(V,L)=H^0(\mathbb P^m,\mathcal O(1))
\otimes H^1(C,K_C)$ and the $i$th coefficient
of this contraction is the perfect Serre-duality
pairing with $\delta(z_iH)$. Its rank is exactly
three for every $m\ge2$.

\emph{3. The actual Poincare-family obstruction.}
We record the precise first-order calculation.
For a smooth variety $V$, a square-zero thickening
by $L$ with class $\kappa$, and a line bundle
$\mathcal Q$ on $V\times P$, its obstruction to
extension to $V_2\times P$ is
\begin{equation}\label{eq:C3-line-obstruction}
        \kappa\mathbin{\lrcorner}c_1^{V}(\mathcal Q)
                   \in H^2(V\times P,\operatorname{pr}_V^*L),
\end{equation}
where $c_1^{V}$ is the component of its Atiyah
class in the $V$-differentials. This is
\citet[Theorem 2.6 and formula (2.7.2),
pp.~166--167]{D_Illusie1972Torsors}, applied to
$\mathbb G_m$-torsors. In this line-bundle case
it follows directly from transition functions:
local splittings of the thickening differ by
derivations $d_{ij}$ with values in $L$; on triple
intersections the product of lifted transition
units differs from one by the evaluation of
$d_{ij}$ on $d\log g_{jk}$. The resulting Cech
two-cocycle is the cup product in
\eqref{eq:C3-line-obstruction}. Thus the formula
uses the actual extension obstruction, not just
a tangent-space dimension.

Let $J=\operatorname{Pic}^0(C)$ and let $P$ be an
abelian variety. If a line bundle on $V\times P$
defines a morphism $\psi:P\to J$ carrying $0$ to
$0$, its $H^0(C,K_C)\otimes H^1(P,\mathcal O_P)$
mixed Atiyah component is that of the normalized
Poincare bundle pulled back along $\psi$.
The corresponding linear map is
\[
                  \psi^*:H^1(J,\mathcal O_J)
                                  \longrightarrow H^1(P,\mathcal O_P).
\]
One can verify the normalization by pulling the
Poincare bundle back along the Abel map
$C\to J$: on $C\times C$ it is
$\mathcal O(\Delta-C\times c_0-c_0\times C)$,
whose mixed diagonal class gives the identity
under Serre duality. The Abel map induces an
isomorphism on $H^1(\mathcal O)$, so this
determines the asserted component. The Poincare
universal property and this cohomological
isomorphism are recorded in
\citet[Theorem 1.2 and Lemma 9.5,
pp.~3 and 31--32]{D_MilneJacobians2021}.

For our $V_2$, the component of
\eqref{eq:C3-line-obstruction} in
\[
                H^1(V,L)\otimes H^1(P,\mathcal O_P)
\]
is obtained by applying $c_\kappa$ to the first
factor of this Poincare tensor. Since
$c_\kappa$ is injective, vanishing of the
obstruction forces $\psi^*=0$. The morphism
$\psi$ is a homomorphism by rigidity of abelian
varieties. Over $\mathbb C$, a nonzero
homomorphism of abelian varieties induces a
nonzero map on $H^1(\mathcal O)$: complex
conjugation gives the corresponding map on
holomorphic one-forms, and zero differential
everywhere forces a homomorphism to be constant.
Consequently an actual extension to
$V_2\times P$ forces $\psi=0$.

\emph{4. Ordinary algebraic families on $X$ vanish.}
Take an ordinarily algebraically trivial line
bundle $M$ on $X$. The finite-descent and smooth
compactification proof in C.3j supplies an
abelian variety $P$, a point $a\in P(\mathbb C)$,
and an actual line bundle $\mathcal M$ on
$X\times P$ such that
\[
                    \mathcal M_0\simeq\mathcal O_X,
                    \qquad \mathcal M_a\simeq M.
\]
For a sum of relations one takes a product of
the Picard varieties and the tensor product
of their pulled-back Poincare bundles. This
uses actual Poincare bundles on those source
Picard varieties, not on their possibly
non-fine abelian images.

The scheme $\mathcal Z\times P$ is regular:
$\mathcal Z$ is regular and its projection
from this product is smooth. Extend
$\mathcal M$ to a line bundle
$\widetilde{\mathcal M}$ on this product by
closing the divisor of a rational section.
The product is noetherian and locally factorial,
so the closure is Cartier. Its fibre at $0$
is trivial outside the irreducible divisor
$V$, hence is $\mathcal O_{\mathcal Z}(bV)$
for an integer $b$. Tensoring by the pullback
of its inverse normalizes that fibre to
$\mathcal O_{\mathcal Z}$, without changing
the original family on $X$.

Restrict to $V\times P$. Since the fibre at
$0$ is trivial and $P$ is connected, its
classifying morphism takes values in
$\operatorname{Pic}^0(V)=J$. It extends to
$V_2\times P$, so step 3 shows this morphism
is zero. By the Poincare universal property
or the see-saw principle the restriction is
therefore pulled back from a line bundle on
$P$. Twist by that pullback. This changes
none of the individual fibre isomorphism
classes and makes the restriction to
$V\times P$ trivial.

The unit sequence for the square-zero
thickening now shows that the class on
$V_2\times P$ comes from
\[
\begin{aligned}
 H^1(V\times P,\operatorname{pr}_V^*L)
  ={}&H^1(V,L)\otimes H^0(P,\mathcal O_P)\\
    &\oplus H^0(V,L)\otimes H^1(P,\mathcal O_P).
\end{aligned}
\]
Restriction to a point of $P$ kills the second
summand and is the same map on the first
summand for every point, since
$H^0(P,\mathcal O_P)=\mathbb C$. Thus all
its individual fibres have the same class
in $\operatorname{Pic}(V_2)$. The fibre at
$0$ is trivial, so the fibre at $a$ is
trivial as well.

For every $r\ge2$,
\[
                  H^1(V,L^r)=H^2(V,L^r)=0.
\]
This follows from the projective-space
cohomology and $H^1(C,K_C^r)=0$ for $r\ge2$.
If $V_r$ is defined by $\mathcal I_V^r$,
the exact unit sequences therefore give
isomorphisms
$\operatorname{Pic}(V_{r+1})\to\operatorname{Pic}(V_r)$
for $r\ge2$. They also make the transition
maps on global units surjective. It follows
inductively that $\widetilde{\mathcal M}_a$
admits compatible trivializations on every
$V_r$, and is formally trivial along $V$.

We spell out why this proves triviality on
the strict local puncture. On $\mathcal Z$,
\[
                        \mathfrak m_R\mathcal O_{\mathcal Z}
                                          =\mathcal I_V.
\]
The old cone coordinates generate the ideal
$(t)$ on every Segre chart; every additional
function from normalization vanishes on the
reduced fibre and hence also belongs to
$(t)$. These two inclusions give the equality.
After completing $R$, full faithfulness of
formal completion for a proper scheme
\citep[Lemma 30.24.1, Tag 0883]{D_StacksFormalPicard}
turns the compatible formal trivializations
into an isomorphism with the trivial line
bundle on $\mathcal Z_{\widehat R}$. In
particular $M$ becomes trivial on the
puncture of $\operatorname{Spec}\widehat R$.

Finally let $Q=j_*M$, the finite rank-one
reflexive $R$-module extending $M$ across
the closed point. Its local cohomology
$H^0_{\mathfrak m}(Q)$ and
$H^1_{\mathfrak m}(Q)$ vanish. Flat base
change of these groups, and the fact that
the closed fibre of completion is a point,
show that $Q\otimes_R\widehat R$ is the
extension of the pulled-back line bundle
from the completed puncture. It is therefore
$\widehat R$. Faithfully flat descent of
finite-presentation flatness gives that $Q$
is a projective rank-one $R$-module, hence
is free since $R$ is local. Thus $M$ is
trivial. This proves the first equality in
\eqref{eq:C3-ordinary-nonsurjection} without
inferring algebraization from tangent data.

\emph{5. Actual nontrivial target families.}
The projection $\mathcal D\to H_0\to C$
pulls back the Poincare bundle on $C\times J$
to an actual family on $Y\times J$.
It induces an injection
\[
                     \operatorname{Pic}^0(C)(\mathbb C)
                              \longrightarrow\operatorname{Pic}(Y).
\]
Indeed if the pullback of a line bundle $Q_C$
is trivial on $Y$, its extension to the
regular $\mathcal D$ is
$\mathcal O_{\mathcal D}(bH_0)$ for some
integer $b$. Restricting to $H_0$ gives
\[
                    \pi_C^*Q_C\simeq L_H^{-b}.
\]
Its degree on a line of the fibre
$\mathbb P^{m-1}$ is zero on the left and
$-b$ on the right. Thus $b=0$, and the
projective-bundle formula for the Picard
group gives $Q_C=\mathcal O_C$. The image
is ordinarily algebraically trivial:
points of the connected projective variety
$J$ can be joined by a chain of smooth
curves, to which the displayed family can
be restricted.

For an explicit nontorsion witness choose
$u\in\mathbb C$ transcendental over
$\overline{\mathbb Q}$, choose $v$ with
$v^4=-1-u^4$, and choose
$\zeta\in\overline{\mathbb Q}$ with
$\zeta^4=-1$. The points
\[
                         p=(u:v:1),\qquad c_0=(1:\zeta:0)
\]
lie on $C$. The line bundle
$\mathcal O_C(p-c_0)$ has infinite order.
Otherwise its Abel--Jacobi point would be
a torsion point of the Jacobian defined
over $\overline{\mathbb Q}$, hence itself
defined over that field. The Abel map is
a closed immersion, also defined over
$\overline{\mathbb Q}$
\citep[Proposition 2.3, p.~7]{D_MilneJacobians2021},
forcing $p$ to be
algebraic, contrary to the choice of $u$.
Its target pullback is therefore a nonzero
element of $J^1(Y)$, while $J^1(X)=0$.

The finite polynomial and Jacobian-ring
certificates are supplied separately. The
obstruction, proper-family and completion
arguments above, not those finite checks,
establish the actual counterexample.
